% Place this refinement after the existing three-point reconstruction proof.
\subsection{Two points for triangles and a sharp worst case}\label{subsec:two-laplace}

\begin{proposition}[Two arguments for complete triangles]\label{prop:two-laplace}
For every complete $Q\in\C_3$ and any $0<\lambda_1<\lambda_2$, the two full
matrices $\widehat\Psi_Q(\lambda_1),\widehat\Psi_Q(\lambda_2)$ determine all
six rates uniquely within $\C_{\le3}$, without separately calibrating $r,s$.
The inverse is rational and locally Lipschitz at each such source.
\end{proposition}
\begin{proof}
Use $M=(P\widehat\Psi)^{-1}$ and $h=b+d$ from the preceding reconstruction.
The Schur complement gives
\begin{equation}
 M_{xy}(\lambda)=-\frac{ad}{r(\lambda+h)},\qquad
 M_{yx}(\lambda)=-\frac{cb}{s(\lambda+h)}.
 \label{eq:two-offdiagonal}
\end{equation}
Both are nonzero at a complete triangle. Thus
\begin{equation}
 \rho=\frac{M_{xy}(\lambda_1)}{M_{xy}(\lambda_2)}
 =\frac{\lambda_2+h}{\lambda_1+h}>1,
 \qquad h=\frac{\lambda_2-\rho\lambda_1}{\rho-1}.
 \label{eq:two-hidden-escape}
\end{equation}
The two-point formulas for $u_i,v_i$ in the preceding proof now recover
$r,s,a,c$, and Eq.~\eqref{eq:branch} recovers $b,d$. All their denominators
are nonzero at the source. A two-state competitor has zero inverse
off-diagonals; a reciprocal three-state competitor with equal data must
have $ad>0$, which forces its complete support. It is therefore recovered
by the same rational formulas, proving uniqueness in the full stated cap.
\end{proof}

Three points remain necessary in the worst case for uncalibrated rates.
At a hidden leaf attached to $x$, $r,s,a,b>0$ and $c=d=0$, so both
off-diagonals in Eq.~\eqref{eq:two-offdiagonal} vanish. The remaining tests are
$G_x(\lambda)=u+v/(\lambda+b)$, with $u=1/r$, $v=a/r>0$, and $G_y=1/s$.
Given two distinct positive arguments, put
\begin{gather}
 D=\frac{G_x(\lambda_1)-G_x(\lambda_2)}{\lambda_2-\lambda_1}>0,\qquad
 v_\beta=D(\lambda_1+\beta)(\lambda_2+\beta),\\
 u_\beta=G_x(\lambda_1)-D(\lambda_2+\beta).
\end{gather}
Every $\beta>0$ with $u_\beta>0$ yields a leaf generator with
\begin{equation}
 (r,s,a,b,c,d)=(u_\beta^{-1},s,v_\beta/u_\beta,\beta,0,0)
 \label{eq:two-leaf-family}
\end{equation}
having the same two full matrices. An open interval around the source
$\beta=b$ is admissible, and $u_\beta'=-D$ makes the generators distinct.
The $y$-leaf is symmetric. Together with Theorem~\ref{thm:inverse}, this
establishes a worst-case minimum of three positive arguments. In the present
reciprocal three-state model with the observed $x$--$y$ pair present, the two
leaf attachments are the only noncomplete irreducible supports; the conclusion
is not a classification of arbitrary sparse larger networks.

For an explicit check, using the tuple order \emph{$(r,s,a,b,c,d)$}, take
\begin{equation}
 \theta_1=(1,1,1,1,0,0),\qquad
 \theta_2=(6/5,1,12/5,2,0,0).
\end{equation}
Direct evaluation of $R(\lambda I-T)^{-1}B$, including the two observed-rate
amplitudes and the $(+,-)$ row order, gives
\begin{equation}
 \widehat\Psi_1(1)=\widehat\Psi_2(1)=
 \begin{pmatrix}0&1/2\\2/5&0\end{pmatrix},\qquad
 \widehat\Psi_1(2)=\widehat\Psi_2(2)=
 \begin{pmatrix}0&1/3\\3/11&0\end{pmatrix}.
 \label{eq:two-leaf-counterexample}
\end{equation}
The complete transforms differ, since
\begin{equation}
 [\widehat\Psi_1-\widehat\Psi_2]_{-,+}(\lambda)
 =-\frac{\lambda(\lambda-1)(\lambda-2)}
 {(\lambda^2+3\lambda+1)(5\lambda^2+28\lambda+12)}.
\end{equation}
At $\lambda=3$ their lower-left entries are $4/19$ and $10/47$.
These are reciprocal trees and both have zero entropy. The example is
ambiguity of the generator from a finite set of summaries, not equality of
the full waiting laws or a nontrivial entropy fiber.

Calibration changes this comparison. The observed rate $r$ is $1$ versus
$6/5$, whereas $s=1$ in both. Their stationary mark intensities in physical
time are $1/3$ per mark versus $6/17$ per mark; their normalized mark
proportions are nevertheless $(1/2,1/2)$ in both. Their traces also differ.
Thus fixed observed rates, a supplied absolute stationary event intensity,
or a supplied trace exclude this particular pair. More generally, if both
$r,s$ are known, subtracting the appropriate constant from $G_i$ lets the
ratio of two nonzero residual tests recover $h$ even at a leaf. Two positive
arguments then suffice for every three-state source in that calibrated class.

One argument is insufficient in general even at complete support: in tuple
order $(r,s,a,b,c,d)$, $(1,1,1,1,1,1)$ and $(2,2,3,2/5,3,2/5)$ share at
$\lambda=1$ the full matrix
$\bigl(\begin{smallmatrix}1/21&8/21\\8/21&1/21\end{smallmatrix}\bigr)$
but differ at $\lambda=2$. Exact symbolic checks of these statements use the
full three-state generators. They also give Jacobian ranks six for two
points at a triangle, five for two points at the displayed leaf, and six
for three points there. Thus the two-point refinement supplies no uniform
raw-rate conditioning guarantee near leaf boundaries; the three-point
stability result retains its separate role.

% SPLIT: EMPIRICAL TV AND FINITE-N MOMENT COVERAGE
% Place this part in the finite-record section near the operational estimator.
\subsection{True-law distance and finite-record confidence sets}
\label{subsec:empirical-topology}

The observation topology in the stability theorems is a distance between
\emph{true} row laws. It is not the total-variation estimation error of a raw
empirical measure. For every finite nonempty row,
$\widehat\mu_I=n_I^{-1}\sum_{m:I_m=I}\delta_{(\tau_m,J_m)}$ is atomic, while
the CTMC row law has a density in time. The finite set of observed atoms has
empirical probability one and true probability zero, so
\begin{equation}
 \TV(\widehat\mu_I,\mu_I)=1.\label{eq:empirical-tv-one}
\end{equation}
Weak convergence and convergence of bounded Laplace tests remain possible.
The moment estimator above uses precisely those bounded tests and therefore
does not assume a shrinking empirical-TV error. A smoothing or detector model
would need its own error analysis.

An elementary confidence rectangle gives a finite-sample alternative to
the pointwise delta approximation.
\begin{proposition}[Sequential moment confidence rectangle]
\label{prop:moment-rectangle}
Fix $K$ positive Laplace arguments in advance, $0<\alpha<1$, and a deterministic
number $n\ge1$ of complete inter-mark intervals. Let $d_{\rm mom}=4K$ and
\begin{equation}
 t_{n,\alpha}=\sqrt{\frac n2\log\frac{2d_{\rm mom}}\alpha}.
\end{equation}
For the empirical joint moments~\eqref{eq:empirical-laplace},
\begin{equation}
 \mathbb P_Q\left\{
 \left|\widehat\mu_{I,jL}-\widehat\Psi_{Q,IL}(\lambda_j)\right|
 \le\frac{t_{n,\alpha}}{n_I}
 \ \text{for all }(I,j,L)\text{ with }n_I>0\right\}\ge1-\alpha.
 \label{eq:finite-moment-rectangle}
\end{equation}
Rows with $n_I=0$ impose no constraint. Stationarity of the initial mark is
unnecessary; the exact CTMC and resolved-reset observation assumptions suffice.
\end{proposition}
\begin{proof}
Let $\mathcal F_0$ contain the first observed mark and let $\mathcal F_m$
include the next $m$ complete marked intervals. Then $I_m$ is
$\mathcal F_{m-1}$-measurable. For one fixed coordinate define
\begin{equation}
 Z_m=\mathbf1\{I_m=I\}
 \left(e^{-\lambda_j\tau_m}\mathbf1\{J_m=L\}
       -\widehat\Psi_{Q,IL}(\lambda_j)\right).
\end{equation}
The reset property gives $\mathbb E[Z_m\mid\mathcal F_{m-1}]=0$.
Conditionally, $Z_m$ lies in an interval of length at most one, so
conditional Hoeffding's lemma yields
$\mathbb E[e^{\eta Z_m}\mid\mathcal F_{m-1}]\le e^{\eta^2/8}$.
Iteration and Chernoff's bound, optimized at $\eta=4t/n$, imply
\begin{equation}
 \mathbb P_Q\left\{\left|\sum_{m=1}^nZ_m\right|>t\right\}
 \le2e^{-2t^2/n}.
\end{equation}
Union over $d_{\rm mom}$ coordinates and use
$\sum_mZ_m=n_I(\widehat\mu_{I,jL}-\widehat\Psi_{Q,IL}(\lambda_j))$
for nonempty rows. The argument neither assumes independent successive rows
nor conditions on their random counts.
\end{proof}

Intersecting the simultaneous rectangle with a declared microscopic model
class gives a confidence set for the generator with coverage at least
$1-\alpha$. Its full entropy image, or a certified enclosing interval,
inherits this coverage and may have an infinite upper endpoint. Finding a
large feasible entropy by numerical optimization does not certify that
endpoint from above. The bound is conservative and does not establish
efficient inversion or an efficient estimator. It applies to fixed $n$
complete intervals and preselected arguments; an adaptive stopping rule,
same-record design selection, terminal censoring at fixed physical time,
or missed events requires additional analysis.

The existing simulations already use actual sequential CTMC records and
dependence-aware moment covariance. Their reported bias, variance, and
Monte Carlo coverage remain operationally meaningful. Their Wald intervals
and low-count warning retain the stated pointwise or heuristic status;
Eq.~\eqref{eq:finite-moment-rectangle} is a separate finite-sample
confidence-set construction, not a retroactive coverage guarantee for those
intervals.
